\part[3]\hfill

  Consider the following sequence of declarations, evaluated in order:

  \begin{codeblock}
    val a = 1
    val f = fn () => 1 div 0
    val g = fn x => raise Fail "boom"
    val res = a + a
  \end{codeblock}

  Does evaluating these declarations raise an exception? Answer yes or no, and
  justify your answer in one sentence. If it does not raise, state the value of
  \code{res}.

  \begin{solutionorbox}[9em]
    No. Binding \code{f} and \code{g} does \textbf{not} run their bodies --- a
    function's body is only evaluated when the function is \textit{applied},
    and neither \code{f} nor \code{g} is ever applied. So no division by zero
    and no \code{raise} occurs, and \code{res} $\eeq$ \code{2}.

    (Contrast \code{val c = 1 div 0}, which is not a function and \textit{would}
    raise immediately, because its right-hand side is evaluated at binding
    time.)
  \end{solutionorbox}


% -----------------------------------------------------------------------------
% 2. Static scope: a closure captures the binding in effect at its definition,
%    not at its call. Misunderstanding: students think name lookup happens at
%    call time against the "current" binding.
% -----------------------------------------------------------------------------
